\documentclass[11pt,a4paper]{article}
\usepackage[italian]{babel}
\usepackage[margin=2.5cm]{geometry}
\usepackage{parskip}
\usepackage{amsmath, amssymb, amsthm}
\usepackage{booktabs}
\usepackage{array}
\usepackage{xcolor}
\usepackage{needspace}
\usepackage{enumitem}
\usepackage{listings}
\usepackage{tikz}
\usepackage[most]{tcolorbox}
\usepackage[hidelinks]{hyperref}
\usetikzlibrary{decorations.pathreplacing, shapes.geometric, arrows.meta, positioning}

% Niente vedove/orfane: i paragrafi non lasciano righe isolate a fine/inizio pagina
\widowpenalty=10000
\clubpenalty=10000

% Elenchi più compatti
\setlist{itemsep=2pt, parsep=0pt, topsep=4pt, partopsep=0pt}

% --- Riquadri con bordo nero, senza numero (si spezzano tra due pagine) ---
% Uso: \begin{definizione} ... \end{definizione}
%      \begin{teorema}[Nome] ... \end{teorema}  (il nome è facoltativo)
\tcbset{nota/.style={enhanced jigsaw, breakable, colback=white,
  colframe=black, boxrule=0.5pt, sharp corners}}
\NewTColorBox{definizione}{}{nota,
  before upper={\textbf{Definizione.}\ }}
\NewTColorBox{teorema}{o}{nota,
  before upper={\textbf{Teorema\IfValueT{#1}{ (#1)}.}\ }}
\NewTColorBox{proposizione}{o}{nota,
  before upper={\textbf{Proposizione\IfValueT{#1}{ (#1)}.}\ }}
\NewTColorBox{proprieta}{o}{nota,
  before upper={\textbf{Proprietà\IfValueT{#1}{ (#1)}.}\ }}
\NewTColorBox{assioma}{o}{nota, boxrule=1pt,
  before upper={\textbf{Assioma\IfValueT{#1}{ (#1)}.}\ }}

% --- Ambienti semplici (senza riquadro) ---
% Il titolo sta in cima al blocco, anche se il contenuto inizia con un riquadro o
% con codice; e non resta mai da solo in fondo a una pagina.
% Uso: \begin{esempio}[Titolo facoltativo] ... \end{esempio}
\newcommand{\newrunin}[2]{%
  \NewDocumentEnvironment{#1}{o}{%
    \par\Needspace{4\baselineskip}\addvspace{\medskipamount}\noindent
    \textbf{#2}\IfValueT{##1}{ (##1)}\textbf{.}\ \ignorespaces}%
  {\par\addvspace{\medskipamount}}}
\newrunin{esempio}{Esempio}
\newrunin{esercizio}{Esercizio}
\newrunin{osservazione}{Osservazione}

% V e F nelle tavole di verità
\newcommand{\V}{\textsf{V}}
\newcommand{\F}{\textsf{F}}

% --- Codice C ---
% Uso: \begin{codice} ... \end{codice}      codice C numerato
%      \begin{comando} ... \end{comando}    comandi da terminale
%      \begin{risultato} ... \end{risultato} output di un programma
%      \lstinline|printf("ciao");|           codice dentro al testo
\lstdefinestyle{c}{
  language=C,
  basicstyle=\ttfamily\small,
  keywordstyle=\color{blue}\bfseries,
  commentstyle=\color{gray}\itshape,
  stringstyle=\color{red!70!black},
  numbers=left,
  numberstyle=\tiny\color{gray},
  numbersep=8pt,
  columns=flexible,
  keepspaces=true,
  breaklines=true,
  showstringspaces=false,
  tabsize=4,
  morekeywords={include, define},
  % lettere accentate nei commenti
  literate={à}{{\`a}}1 {è}{{\`e}}1 {é}{{\'e}}1 {ì}{{\`i}}1
           {ò}{{\`o}}1 {ù}{{\`u}}1 {È}{{\`E}}1 {À}{{\`A}}1
}
\lstset{style=c}
\newtcblisting{codice}{listing only, nota, left=6mm,
  listing options={style=c}}
\newtcblisting{comando}{listing only, nota,
  listing options={basicstyle=\ttfamily\small, numbers=none, language={},
    columns=flexible, keepspaces=true, breaklines=true}}
\newtcblisting{risultato}{listing only, enhanced jigsaw, breakable,
  colback=black!5, colframe=black, boxrule=0.5pt, sharp corners,
  listing options={basicstyle=\ttfamily\small, numbers=none, language={},
    columns=flexible, keepspaces=true, breaklines=true}}

% --- Diagrammi di flusso (simboli standard) ---
\tikzset{
  terminale/.style = {ellipse, draw, minimum width=2.5cm, minimum height=0.9cm},
  io/.style        = {trapezium, trapezium left angle=70, trapezium right angle=110,
                      draw, minimum width=2.5cm, minimum height=0.9cm},
  processo/.style  = {rectangle, draw, minimum width=2.5cm, minimum height=0.9cm},
  decisione/.style = {diamond, draw, aspect=2, minimum width=2.5cm},
  freccia/.style   = {thick, -{Stealth}}
}

\title{Funzioni elementari}
\author{Marco Cerbella}
\date{Lezione 6 -- 12 ottobre 2026}

\begin{document}
\maketitle

\section{Trasformazioni geometriche dei grafici}

Data $f : \mathbb{R} \to \mathbb{R}$, il \emph{valore assoluto} di $f$ è
\[
|f(x)| =
\begin{cases}
f(x) & \text{se } f(x) \ge 0, \\
-f(x) & \text{se } f(x) < 0.
\end{cases}
\]
Con semplici trasformazioni geometriche si disegna il grafico delle funzioni seguenti a partire da quello di $f$.

\begin{center}
\renewcommand{\arraystretch}{1.3}
\begin{tabular}{@{}lp{9.5cm}@{}}
\toprule
Funzione & Effetto sul grafico di $f$ \\
\midrule
$y_1 = f(x) + a$, $a \in \mathbb{R}$ & traslazione in verticale: verso l'alto se $a > 0$, verso il basso se $a < 0$ \\
$y_2 = f(x + a)$, $a \in \mathbb{R}$ & traslazione in orizzontale: verso sinistra se $a > 0$, verso destra se $a < 0$ \\
$y_3 = k \cdot f(x)$, $k > 0$ & dilatazione in verticale se $k > 1$ (il grafico si allontana dall'asse $x$), contrazione se $0 < k < 1$ (si schiaccia sull'asse $x$) \\
$y_4 = f(kx)$, $k > 0$ & dilatazione in orizzontale se $0 < k < 1$ (si allontana dall'asse $y$), contrazione se $k > 1$ (si schiaccia sull'asse $y$) \\
$y_5 = |f(x)|$ & la parte di grafico con $f(x) < 0$ viene riflessa rispetto all'asse $x$ \\
$y_6 = f(|x|)$ & si tiene la parte con $x \ge 0$ e la si riflette rispetto all'asse $y$ (la parte con $x < 0$ viene sostituita da questa riflessione) \\
\bottomrule
\end{tabular}
\end{center}

\begin{esempio}[Parabola]
Sia $f(x) = x^2$. Allora $g(x) = x^2 + 1$ è la parabola traslata verso l'alto di $1$ e $h(x) = x^2 - 1$ quella traslata verso il basso di $1$. La funzione $g(x) = (x+1)^2$ si annulla per $x = -1$ (traslazione a sinistra di $1$) e $h(x) = (x-1)^2$ si annulla per $x = 1$ (traslazione a destra), mentre $f(0) = 0$.
\end{esempio}

\begin{esempio}[Seno]
Per $x \in [0, 2\pi]$, sia $f(x) = \sin x$.
\begin{itemize}
    \item $g(x) = 2\sin x$ e $h(x) = \tfrac{1}{2}\sin x$ hanno gli stessi zeri di $f$ ($0, \pi, 2\pi$) ma massimo $2$ e $\tfrac12$ (in $x = \pi/2$), invece di $1$.
    \item $g(x) = \sin(2x)$ si contrae in orizzontale: si annulla per $x = 0, \tfrac{\pi}{2}, \pi, \tfrac{3\pi}{2}, 2\pi$ e vale $1$ in $x = \tfrac{\pi}{4}$.
    \item $h(x) = \sin\!\left(\tfrac{1}{2}x\right)$ si dilata in orizzontale: si annulla per $x = 0$ e $x = 2\pi$ e vale $1$ in $x = \pi$.
\end{itemize}
\end{esempio}

\begin{osservazione}
Per $y_5 = |f(x)|$ e $y_6 = f(|x|)$ la lezione mostra graficamente due casi: per $|f(x)|$ si ribalta sopra l'asse $x$ la parte di $f$ che sta sotto; per $f(|x|)$ si usa $f(|x|) = f(x)$ se $x \ge 0$ e $f(|x|) = f(-x)$ se $x < 0$.
\end{osservazione}

\section{Funzioni lineari e lineari affini}

Si usa una \emph{funzione lineare} per esprimere una proporzionalità diretta tra due grandezze $x$ e $y$: al raddoppiare di $x$ raddoppia $y$, al triplicare di $x$ triplica $y$.

La legge di una funzione lineare è $y = mx$, quella di una funzione lineare affine è $y = mx + q$, con $m, q \in \mathbb{R}$. Il dominio è $\mathbb{R}$.

\begin{center}
\renewcommand{\arraystretch}{1.3}
\begin{tabular}{@{}lccc@{}}
\toprule
 & $m > 0$ & $m = 0$ & $m < 0$ \\
\midrule
Monotonia & str.\ crescente & costante & str.\ decrescente \\
Limitatezza & non limitata & limitata & non limitata \\
$q = 0$ (lineare) & dispari & pari & dispari \\
$q \ne 0$ (affine) & non simmetrica & pari & non simmetrica \\
\bottomrule
\end{tabular}
\end{center}

Per $m \ne 0$ la funzione è illimitata, strettamente monotona e quindi invertibile. Per $q = 0$ e $m \neq 0$ la funzione è dispari, perché $f(-x) = m(-x) = -mx = -f(x)$. Per $q \ne 0$ non è simmetrica: $f(-x) = -mx + q$ e $-f(x) = -mx - q$ sono diversi.

\section{Funzione valore assoluto}

Il valore assoluto (o modulo) di $x$ è
\[
f(x) = |x| =
\begin{cases}
x & \text{se } x \ge 0, \\
-x & \text{se } x < 0.
\end{cases}
\]
Si ha $|x| \ge 0$ per ogni $x \in \mathbb{R}$ e $|x| = 0 \iff x = 0$. La funzione è \emph{pari} ($f(-x) = |-x| = |x| = f(x)$), limitata inferiormente ma non superiormente, e \emph{non monotona}.

\section{Proporzionalità inversa}

L'altro tipo di proporzionalità tra due grandezze è la \emph{proporzionalità inversa}: $xy = k$, con $k \in \mathbb{R}$. Il prodotto delle due grandezze rimane costante al variare dell'una o dell'altra. Per $k = 1$ si ottiene
\[
f : \mathbb{R} \setminus \{0\} \to \mathbb{R}, \qquad f(x) = \frac{1}{x}.
\]
È illimitata, decrescente per $x > 0$ e per $x < 0$ (non su tutto il dominio) e \emph{dispari}, perché $f(-x) = \dfrac{1}{-x} = -\dfrac{1}{x} = -f(x)$.

\section{Funzioni potenza}

Le funzioni potenza sono del tipo
\[
f(x) = k\,x^{\alpha}, \qquad \alpha \ne 0,\ k, \alpha \in \mathbb{R}.
\]
Il valore $f(x)$ è proporzionale, con costante di proporzionalità $k$, a $x^\alpha$. Si studia $x^\alpha$ al variare di $\alpha$.

\subsection*{Esponente naturale $\alpha = n \in \mathbb{N}$}

\begin{center}
\renewcommand{\arraystretch}{1.3}
\begin{tabular}{@{}lll@{}}
\toprule
 & $n$ pari (es.\ $x^2, x^4$) & $n$ dispari (es.\ $x^3, x^5$) \\
\midrule
Dominio & $\mathbb{R}$ & $\mathbb{R}$ \\
Immagine & $[0,+\infty)$ & $\mathbb{R}$ \\
Simmetria & pari & dispari \\
Invertibile & no & sì \\
\bottomrule
\end{tabular}
\end{center}
Per $n$ dispari la funzione è anche illimitata.

\subsection*{Esponente razionale positivo $\alpha = \frac{m}{n}$, $m, n \in \mathbb{N}$}

\begin{center}
\renewcommand{\arraystretch}{1.3}
\begin{tabular}{@{}llll@{}}
\toprule
 & $n$ pari & $n$ dispari, $m$ dispari & $n$ dispari, $m$ pari \\
\midrule
Esempi & $x^{1/2}, x^{3/2}, x^{5/2}$ & $x^{1/3}, x^{5/3}$ & $x^{2/3}$ \\
Dominio & $[0,+\infty)$ & $\mathbb{R}$ & $\mathbb{R}$ \\
Immagine & $[0,+\infty)$ & $\mathbb{R}$ & $[0,+\infty)$ \\
Simmetria & dominio non simmetrico & dispari & pari \\
Invertibile & su $[0,+\infty)$ & sì & no \\
\bottomrule
\end{tabular}
\end{center}
Quando il dominio non è simmetrico non ha senso dire pari o dispari.

\subsection*{Esponente razionale negativo $\alpha = -\frac{m}{n}$, $m, n \in \mathbb{N}$}

\begin{center}
\renewcommand{\arraystretch}{1.3}
\begin{tabular}{@{}llll@{}}
\toprule
 & $n$ pari & $n$ dispari, $m$ dispari & $n$ dispari, $m$ pari \\
\midrule
Dominio & $]0,+\infty)$ & $\mathbb{R} \setminus \{0\}$ & $\mathbb{R} \setminus \{0\}$ \\
Immagine & $]0,+\infty)$ & $\mathbb{R} \setminus \{0\}$ & $]0,+\infty)$ \\
Simmetria & dominio non simmetrico & dispari & pari \\
Invertibile & su $]0,+\infty)$ & su $\mathbb{R} \setminus \{0\}$ & no \\
\bottomrule
\end{tabular}
\end{center}

\subsection*{Esponente irrazionale $\alpha \in \mathbb{R} \setminus \mathbb{Q}$}

Se si prende come dominio $[0,+\infty)$ oppure $]0,+\infty)$ si possono considerare anche potenze con esponente non razionale.
\begin{itemize}
    \item $\alpha > 0$: dominio $[0,+\infty)$, immagine $[0,+\infty)$, dominio non simmetrico, invertibile su $[0,+\infty)$.
    \item $\alpha < 0$: dominio $]0,+\infty)$, immagine $]0,+\infty)$, dominio non simmetrico. La funzione è strettamente decrescente, quindi invertibile (le slide del corso riportano per errore ``non invertibile'').
\end{itemize}

\subsection*{Riepilogo del dominio}

$x^\alpha$ è definita per:
\begin{itemize}
    \item $x \in \mathbb{R}$ se $\alpha = \frac{m}{n}$ con $m, n \in \mathbb{N}$ e $n$ dispari;
    \item $x \ge 0$ se $\alpha = \frac{m}{n}$ con $m, n \in \mathbb{N}$ e $n$ pari;
    \item $x \ne 0$ se $\alpha = -\frac{m}{n}$ con $m, n \in \mathbb{N}$ e $n$ dispari;
    \item $x > 0$ se $\alpha = -\frac{m}{n}$ con $m, n \in \mathbb{N}$ e $n$ pari;
    \item $x \ge 0$ se $\alpha \in \mathbb{R} \setminus \mathbb{Q}$ e $\alpha > 0$;
    \item $x > 0$ se $\alpha \in \mathbb{R} \setminus \mathbb{Q}$ e $\alpha < 0$.
\end{itemize}

\section{Funzioni polinomiali e razionali}

Una \emph{funzione polinomiale} (o \emph{polinomio}) è una funzione del tipo
\[
P(x) = a_n x^n + \dots + a_1 x + a_0, \qquad a_i \in \mathbb{R},\ a_n \ne 0.
\]
Il \emph{grado} del polinomio è l'esponente della potenza più alta, qui $n$. Il dominio è $\mathbb{R}$.
\begin{itemize}
    \item $P$ è pari se e solo se tutti i coefficienti di indice dispari sono nulli (per esempio $x^4 + x^2$);
    \item $P$ è dispari se e solo se tutti i coefficienti di indice pari sono nulli.
\end{itemize}

Una \emph{funzione razionale} è una funzione del tipo $\dfrac{P(x)}{Q(x)}$ con $P, Q$ polinomi. Il dominio è $\mathbb{R}$ privato degli zeri del denominatore.

\begin{esempio}[Test]
$f(x) = x^2 - 1$ è una funzione: (1) monotona, (2) invertibile, (3) pari, (4) limitata? La risposta corretta è la (3): $f(-x) = (-x)^2 - 1 = f(x)$, mentre $f$ non è monotona, non è iniettiva ($f(1) = f(-1)$) e non è limitata superiormente.
\end{esempio}

\end{document}
